\section{随机过程}

\subsection{实随机过程}
我们前面介绍了多元随机变量，注意到我们之前研究的都是一族有限的随机变量。对多元随机变量\(\mathbf{X}=(X_1,\ldots,X_n)\)来说，指标集是有限集\(T=\{1,2, \ldots, n\}\)。我们可以把它描述为从样本空间\(\Omega\)到\(\mathbb{R}^n\)的一个函数。当我们把指标集\(T\)推广为无限集时，事情就复杂很多了。我们现在面对的是一族无限的随机变量\(\{X_t:t\in T\}\)，我们称这一族随机变量为一个\textbf{随机过程}。借助选择公理，我们可以定义这些随机变量（本质是一族函数）的无穷积——一个随机过程就是一个从\(\Omega\)到\(\mathbb{R}\)的无穷笛卡尔积\(\mathbb{R}^T\)的函数。

一般来说，我们会用实数（正实数、实数区间）、整数（正整数）作为指标集。这就带来了\textbf{连续参数随机过程}和\textbf{离散参数随机过程}的区别。我们一般把连续参数随机过程记作\(X(t)\)，离散参数随机过程记作\(X_n\)。

这带来了一个本质的区别：我们无法再像有限维那样写出一个包含所有变量的“终极”分布函数，我们只能从有限处把握它。我们可以了解其中的任意有限多个随机变量的联合分布。另一方面，我们刚才将随机过程视为\(T\to \mathbb{R}^{\Omega}\)的函数，其中\(\mathbb{R}^{\Omega}\)表示\(\Omega \to \mathbb{R}\)的函数，也就是随机变量，不太正式地说，就是个\(T\to(\Omega \to \mathbb{R})\)的函数。通过反柯里化和柯里化，它自然地对应于一个\((T\times \Omega) \to \mathbb{R}\)这样的二元函数，也就自然对应于一个\(\Omega\to(T \to \mathbb{R})\)的函数，即一个取值为函数的随机变量！某个样本对应的\(T \to \mathbb{R}\)的函数就叫\textbf{样本函数}。

不过从随机函数的这个观点来看，想要这些样本函数有什么好的性质那可就是奢望了。它们往往支离破碎，不可理喻。不过我们也许可以让它\textbf{几乎处处}是“好”的，我们会在后面提到它。但是我们先从另一个视角来看看它，这里有一个更迫切的问题。

前面说了那么多，可是我们如果了解一族无限的随机变量\(\{X_t:t\in T\}\)中任意有限多个随机变量的联合分布，是否真的能自洽地给出一个随机过程，使得使得这些有限维分布恰好就是我们所指定的？这个问题看似简单，实则深刻：我们凭什么相信，这些分散在有限维度上的信息能够拼凑出一个在无限维度上一致的整体？如果它们不能，那我们关于随机过程的一切讨论都将失去根基。

Kolmogorov扩张定理回答了这个问题。它告诉我们：只要这些有限维分布满足两个极其自然的相容性条件，就必然存在一个概率空间和一个定义在其上的随机过程，使得该过程的有限维分布与我们指定的完全一致。这个定理为随机过程的存在性提供了坚实的逻辑基础。

\begin{theorem}[Kolmogorov扩张定理]
设 $T$ 是一个集合。对每个有限子集 ${t_1, t_2, \dots, t_n} \subset T$，令 $\mu_{t_1, t_2, \dots, t_n}$ 是 $\mathbb{R}^n$ 上的一个概率测度。假设这些测度满足以下相容性条件：
\begin{itemize}
    \item 重排不变性：对任意排列 $\pi$ 和任意 $B \in \mathcal{B}(\mathbb{R}^n)$，$\mu_{t_{\pi(1)}, \ldots,t_{\pi(n)}}(B) = \mu_{t_1,\ldots,t_n}(\pi^{-1}(B))$
    \item 边缘相容性：对任意 $m < n$ 和任意 $B \in \mathcal{B}(\mathbb{R}^m)$，$\mu_{t_1,\ldots,t_m}(B) = \mu_{t_1,\ldots,t_n}(B\times \mathbb{R}^{n-m})$
\end{itemize}

那么存在一个概率空间 $(\Omega, \mathcal{F}, P)$ 和一个随机过程 ${X_t: t \in T}$ 定义在其上，使得对任意有限子集 ${t_1, \dots, t_n} \subset T$，随机向量 $(X_{t_1}, \dots, X_{t_n})$ 的联合分布恰为 $\mu_{t_1, \dots, t_n}$。
\end{theorem}

前面说完了从相容的有限维分布得出随机过程，下面我们从随机过程给出它有限维的特性。对于随机过程 \(X(t)\)，其有限维分布函数定义为：
对任意正整数 \(n\)，任意 \(t_1, t_2, \ldots, t_n \in T\)，以及任意实数 \(x_1, \ldots, x_n\)，
有
\[
F_{X}(x_1,\ldots,x_n;t_1,\ldots,t_n) = P\left( X(t_1) \le x_1, \ldots, X(t_n) \le x_n \right).
\]
若分布函数关于每个变量可微，则相应的有限维概率密度函数为
\[
f_{X}(x_1,\ldots,x_n;t_1,\ldots,t_n) = \frac{\partial^n F_{t_1,\ldots,t_n}(x_1,\ldots,x_n)}{\partial x_1 \cdots \partial x_n}.
\]
所有这样的有限维分布构成的集合称为随机过程的有限维分布族，它刻画了随机过程的有限维统计特性。

我们接着直接走向它的数字特征。

我们先从均值开始。显然如果我们把随机过程看成一个随机函数的话，那它的均值（如果存在的话）也应该是一个函数。我们可以很自然地定义函数的均值为均值的函数。对于每个固定的指标\(t\in T\)，我们都可以给出\(X_t\)的均值（如果存在）。我们把随机过程\(X(t)\)的均值记作\(\E[X]\)或者\(m_X(t)\)。定义\(m_X(t)=\E[X(t)]\)。

均值可以不存在，一个简单的例子就是对于任意\(t\in T\)，\(X_t\sim Cauchy(0,1)\)的随机过程。不过一般情况下我们都喜欢处理均值存在的随机过程，即\textbf{一阶矩过程}。更进一步，如果对于任意\(t\)，\(\E[X^2(t)]\)都存在，我们就将其称为\textbf{二阶矩过程}。用一点柯西不等式，我们可以知道：
\begin{gather*}
|\E[X(t)]| \leq \E[|X(t)|] \leq \sqrt{\E[X^2(t)]}\\
(\E[X(s)X(t)])^2 \leq \E[X^2(s)]\E[X^2(t)]\\
\end{gather*}
于是二阶矩过程的一阶矩存在，“混合二阶矩”也存在。

我们可以沿着前面的想法继续：研究某些确定的\(s,t\in T\)给出的随机变量\(X(s),X(t)\)，通过它们的数值特征来描述整个随机过程。我们先从协方差入手，对于给定的\(s,t\in T\)我们可以定义\(\Cov(X(s),X(t))\)。这时不妨让它们变成可以改变的参数，我们就得到了\textbf{自协方差函数}：
\[
C_{XX}(s, t) = \Cov(X(s),X(t)),\quad s,t\in T
\]
把\(s\)和\(t\)同一，就得到\textbf{方差函数}：
\[
\sigma^2_{X}(t) = C_{XX}(t, t)
\]
开个方就是\textbf{标准差函数}或者\textbf{均方差函数}：
\[
\sigma_{X}(t) = \sqrt{\sigma^2_{X}(t)}
\]
我们知道协方差是个内积（如果加上\textbf{几乎处处}和零均值的条件），\(\E[XY]\)也是个内积（如果加上\textbf{几乎处处}的条件），但是比较特别地，我们在研究随机过程的时候还真就比较喜欢\(\E[XY]\)这个内积，我们称之为\textbf{自相关函数}：
\[
R_{XX}(s,t)=\E[X(s)X(t)] = C_{XX}(s,t)+m_X(s)m_X(t)
\]
把\(s\)和\(t\)同一，就得到\textbf{均方值函数}：
\[
\psi_X(t)=R_{XX}(t,t)
\]
在不至于混淆的时候，我们也把\(R_{XX}\)和\(C_{XX}\)写成\(R_X\)和\(C_X\)。

同样地，对于两个不同的随机过程\(X(t)\)和\(Y(t)\)，我们也可以定义它们之间相互的\textbf{互相关函数}和\textbf{互协方差函数}：
\begin{gather*}
R_{XY}(s, t) = \E[X(s)Y(t)]\\
C_{XY}(s, t) = \Cov[X(s)Y(t)]
\end{gather*}
\begin{attention}
\(R_{XY}\)和\(R_{YX}\)是不同的，一般来说我们只能知道\(R_{XY}(s,t)=R_{YX}(t,s)\)。\(C_{XY}\)和\(C_{YX}\)也是这样的。
\end{attention}

对随机变量\(X,Y\)来说，我们把\(\Cov(X,Y)=0\)称为\(X,Y\)不相关。对于多维随机变量，它们中的任意维不相关我们才称之为不相关。对随机过程\(X(t),Y(t)\)来说，若对于任意\(s,t\in T\)，\(C_{XY}(s,t)=0\)，我们称\(X(t)\)与\(Y(t)\)不相关。而对于相关函数这种内积，若对于任意\(s,t\in T\)，\(R_{XY}(s,t)=0\)，我们称\(X(t)\)与\(Y(t)\)正交。

这时我们会想到随机变量的独立性和相关性的关系。如果要把它推广到随机过程来，我们必须先定义独立性的概念。和多元随机变量的独立性类似，我们需要考虑所有维度的独立性。但是我们这里不能这么做，因为随机过程没有包含所有维度的无限维分布函数，我们只能这样定义：随机过程\(X(t),Y(t)\)独立当且仅当\(X(t)\)中的任意\(n\)维和\(Y(t)\)中的任意\(m\)维独立。具体地说，对于任意正整数\(n,m\)和任意\(t_1,\ldots,t_n;t'_1,\ldots,t'_m\in T\)，定义\(X(t),Y(t)\)的\(n+m\)维分布函数为：

\begin{align*}
&F(x_1,\ldots,x_n, y_1,\ldots,y_m;t_1,\ldots,t_n; t'_1,\ldots,t'_m)\\
= &P\left( X(t_1) \le x_1, \ldots, X(t_n) \le x_n, Y(t'_1) \le y_1, \ldots, Y(t'_m) \le y_m \right),
\end{align*}

我们把它紧凑地写成\(F(\mathbf{x},\mathbf{y};\mathbf{t},\mathbf{t}')\)。接下来\(X(t)\)和\(Y(t)\)独立当且仅当对于任意正整数\(n,m\)和任意\(\mathbf{t}\in T^n;\mathbf{t}'\in T^m\)和任意\(\mathbf{x}\in\mathbb{R}^n;\mathbf{y}\in\mathbb{R}^m\)
\[
F(\mathbf{x},\mathbf{y};\mathbf{t},\mathbf{t}')=F_X(\mathbf{x};\mathbf{t})F_Y(\mathbf{y};\mathbf{t}')
\]

对于随机过程来说，还有另一种有意义的平均：即对\(t\)取平均，我们一般将其称为\textbf{时间平均}。前一种平均则被称为\textbf{统计平均}。如果我们像前面所说，把一个随机过程视为一个\(\Omega\times T \to \mathbb{R}\)的函数，写成\(X(e,t)\)，那么我们可以说：统计平均是按\(e\)平均，取遍样本空间，得到一个\(T\to \mathbb{R}\)的函数；时间平均是按\(t\)平均，取遍整个\(T\)，得到一个\(\Omega \to \mathbb{R}\)的函数（也就是随机变量）。

为了和均值相区分，我们把\(X(t)\)的时间平均记作\(\overline{X(t)}\)我们这样定义时间平均：
\[
\overline{X(t)} = \lim_{T\to\inf}\left(\frac{1}{T}\int_{-T/2}^{T/2}X(t)\dif t\right)
\]

我们这里用的是\(T=\mathbb{R}\)的情形。不过我们需要对这个式子保持警惕，因为我们还不知道在什么意义上可以对随机变量做积分（这可是随机的！）。同时，这个极限不一定存在。我们稍后给出随机变量的分析。

时间平均和统计平均之间的关系不是那么简单的。它们往往不相等，而它们\textbf{几乎必然}相等则称为均值具有\textbf{遍历性}。遍历性又叫\textbf{各态历经性}，这个名字的由来并不是非常显然。这个词源于玻尔兹曼在统计力学中的遍历假设：一个孤立的力学系统在长时间演化下会历经所有具有相同能量的微观状态，因此时间平均应该等于统计平均（即所有可能微观状态的等概率平均）。后来这一概念被引入随机过程理论，借用了物理上的直觉。它的名字 “ergodic” 来自希腊语 ergon（功） + hodos（路径），原意是“沿着能量面走遍所有点”的那种路径。

\subsection{复随机过程}
到目前为止，我们讨论的都是实随机过程，我们下面开启复随机过程的讨论。事实上复随机过程没有什么特别不同的。它可以由一族无穷的复随机变量给出，也可以通过两个实随机过程给出。它的数字特征和是随机过程基本相同，只是所有内积都得从实内积变成酉内积。我们在这里依然遵守第一个位置线性，第二个位置共轭线性的定则。

\begin{gather*}
R_{XY}(s,t) = \E[X(s)Y(t)]\\
C_{XY}(s,t) = \Cov(X(s),Y(t)) = R_{XY}(s,t)-m_X(s)m_Y^*(t)\\
\end{gather*}

我们之前的对称性也变成了共轭对称性：
\begin{gather*}
R_{XY}(s,t) = R_{YX}^*(t,s)\\
C_{XY}(s,t) = C_{YX}^*(t,s)\\
\end{gather*}

\subsection{随机分析}
上面我们引入了时间平均的概念，但是它显然是可疑的。最大的一个问题是，我们要怎么定义随机变量的积分呢？我们可以先从一个朴素一点的想法入手：对于任意给定的\(e\in \Omega\)，随机过程\(X(e,t)\)都给出一个函数，我们逐点进行积分就好了。最后我们就应该可以得到一个随机变量：对于每个样本，都对应一个积分值。

但这仍然是值得怀疑的：并不是每个样本函数都可以很好地进行积分，我们需要一些正则性条件保证积分都可以进行。我们有一些正则性定理来帮助我们。Kolmogorov-Chentsov正则性定理就告诉我们，只要满足这些条件，就存在一个和\(X(t)\)（几乎必然意义下）等价的\(\tilde{X}(t)\)使得样本函数都局部可积。一般来说这个正则条件是很容易满足的，但仍然并非所有过程都满足。

\begin{theorem}[Kolmogorov-Chentsov 正则性定理]
设 $\{X_t : t \in [0,1]\}$（或更一般地，$t$ 属于 $\mathbb{R}^d$ 中的一个有界区域）是一个随机过程。假设存在常数 $\alpha > 0, \beta > 0, C > 0$，使得对任意 $s, t \in [0,1]$ 都有：
\[
\E[|X_t - X_s|^\alpha] \le C |t - s|^{d + \beta}
\]
其中 $d$ 是指标集的维数（当指标集为 $[0,1]$ 时 $d=1$）。

那么存在一个过程 $\{\tilde{X}_t : t \in [0,1]\}$，满足：

\begin{itemize}
    \item 对每个 $t \in [0,1]$，有 $\tilde{X}_t = X_t$ 几乎必然成立。
    \item 几乎所有样本轨道 $t \mapsto \tilde{X}_t(\omega)$ 是\textbf{连续}的。更进一步，对任意 $\gamma \in (0, \beta/\alpha)$，样本轨道具有 $\gamma$-Hölder 连续性，即存在一个几乎必然有限的随机变量 $K_\gamma(\omega)$ 使得：
    \[
    |\tilde{X}_t(\omega) - \tilde{X}_s(\omega)| \le K_\gamma(\omega) |t-s|^\gamma, \quad \forall s, t \in [0,1]
    \]
\end{itemize}

特别地，当指标集为 $\mathbb{R}^d$ 中的有界区域时，只需将条件中的 $|t-s|$ 替换为欧氏距离 $\|t-s\|$，结论在任意有界闭区域上成立。
\end{theorem}

这时就有几个问题了：不带正则性条件可以吗？积分和极限、期望等算子交换吗？没有这么好。我们需要一致收敛性或者Fubini定理来交换，但这些条件对随机过程来说往往过于苛刻。于是我们可以转向另一种积分：均方积分。它建立在 $L^2(\Omega)$ 空间（也就是均方收敛的随机变量空间）的框架之上，将积分视为均方意义下的极限。我们谈论它需要从均方收敛性开始说起。

\subsubsection{均方收敛}

\begin{definition}[均方收敛]
  设 $\{X_n\}_{n\ge1}$ 和 $X$ 是二阶矩有限的随机变量。称 $X_n$ \textbf{均方收敛}于 $X$，记作
  \[
  X_n \xrightarrow{L^2} X \quad \text{或} \quad \operatorname*{l.i.m.}_{n\to\infty} X_n = X,
  \]
  如果
  \[
  \lim_{n\to\infty} \E\big[|X_n - X|^2\big] = 0.
  \]
\end{definition}

均方收敛具有以下优良性质：

\begin{enumerate}
  \item \textbf{完备性}：$L^2(\Omega)$ 是希尔伯特空间，因此柯西列必收敛。
  \item \textbf{期望连续性}：若 $X_n \xrightarrow{L^2} X$，则 $\E[X_n] \to \E[X]$。
  \item \textbf{内积连续性}：若 $X_n \xrightarrow{L^2} X$，$Y_n \xrightarrow{L^2} Y$，则 $\E[X_n Y_n] \to \E[XY]$。
\end{enumerate}

我们后面在求随机过程的极限时都默认是均方意义下的。
我们可以继续定义均方连续性和均方微分，但是我就在这里略过了。

\subsubsection{均方积分}

利用均方收敛，我们可以定义均方积分。

\begin{definition}[均方积分]
设 $\{X(t), t\in[a,b]\}$ 是随机过程，满足对每个 $t$，$\E[X^2(t)]<\infty$。考虑 $[a,b]$ 的分割
  \[
  \Delta: a = t_0 < t_1 < \cdots < t_n = b,
  \]
  记 $\|\Delta\| = \max_i (t_i - t_{i-1})$。选取 $\xi_i \in [t_{i-1}, t_i]$，构造黎曼和
  \[
  S(\Delta) = \sum_{i=1}^n X(\xi_i) (t_i - t_{i-1}).
  \]
  若存在随机变量 $I$ 使得
  \[
  \lim_{\|\Delta\|\to 0} \E\big[|S(\Delta) - I|^2\big] = 0,
  \]
  则称 $X(t)$ 在 $[a,b]$ 上\textbf{均方可积}，并称 $I$ 为均方积分，记作
  \[
  I = \int_a^b X(t)\,\dif t \quad \text{(均方意义下)}.
  \]
\end{definition}

均方积分的存在性由以下定理保证。

\begin{theorem}[均方可积的充分条件]
  设 $X(t)$ 满足：
  \begin{enumerate}
    \item $X(t)$ 是\textbf{均方可测}的，即映射 $(e, t) \mapsto X(e, t)$ 关于 $\mathcal{F}\otimes\mathcal{B}([a,b])$ 可测；
    \item $\E[X^2(t)]$ 在 $[a,b]$ 上勒贝格可积。
  \end{enumerate}
  则 $X(t)$ 在 $[a,b]$ 上均方可积。进一步，若 $X(t)$ 是\textbf{均方连续}的（即 $\lim_{h\to0} \E[|X(t+h)-X(t)|^2]=0$），则黎曼和在均方意义下收敛到积分值。
\end{theorem}

均方积分具有以下特别的性质。设 $X(t), Y(t)$ 在 $[a,b]$ 上均方可积，$c$ 为常数。则：
\begin{enumerate}
\item \textbf{期望与积分交换}：$\E\left[\int_a^b X(t)\,\dif t\right] = \int_a^b \E[X(t)]\,\dif t$。
\item \textbf{协方差与积分交换}：
\[
\Cov\left(\int_a^b X(s)\,\dif s, \int_a^b Y(t)\,\dif t\right) = \int_a^b\int_a^b \Cov(X(s), Y(t))\,\dif s\,\dif t.
\]
\item \textbf{均方连续性}：$F(T) = \int_a^T X(t)\,\dif t$ 是 $T$ 的均方连续函数。
\end{enumerate}

\begin{attention}
期望与积分的交换在均方框架下\textbf{自动成立}，无需额外条件。这正是均方积分相较于逐点积分的核心优势——它绕开了 Fubini 定理的条件验证，交换性就在定义之中。
\end{attention}
    
\subsection{随机过程的性质}

在随机过程的理论与应用中，我们通常关注其若干核心概率性质，这些性质常常作为模型假设被直接使用。常见的性质包括马尔可夫性、平稳性以及独立增量性。下面分别予以详细阐述。

\subsubsection{马尔可夫性}

马尔可夫性，又称无记忆性，是随机过程中最基本且应用最广泛的性质之一。它刻画了过程的“现在决定未来”的本质：在已知当前状态的条件下，过程的未来演化与过去状态无关。

\begin{definition}[马尔可夫性]
设 \(X(t)\) 为一随机过程，状态空间为 \(\mathcal{S}\)。若对任意 \(n \ge 1\)，任意时刻 \(t_1 < t_2 < \dots < t_n < t\) 以及任意状态 \(x_1, x_2, \dots, x_n, x \in \mathcal{S}\)，有
\[
P(X_t = x \mid X_{t_1} = x_1, \dots, X_{t_n} = x_n) = P(X_t = x \mid X_{t_n} = x_n),
\]
则称该过程具有马尔可夫性（或称为马尔可夫过程）。对于连续时间情形，上述等式在\textbf{几乎必然}意义下成立。
\end{definition}

具有马尔可夫性的过程称为马尔可夫过程。根据状态空间和时间参数的不同，可进一步分为：
\begin{itemize}
    \item \textbf{马尔可夫链}：时间离散、状态离散（或可测）的马尔可夫过程；
    \item \textbf{马尔可夫跳跃过程}：时间连续、状态离散；
    \item \textbf{扩散过程}：时间连续、状态连续。
\end{itemize}

\subsubsection{平稳性}

平稳性刻画了随机过程在时间平移下的不变性。根据要求强弱的不同，分为\textbf{严平稳}和\textbf{宽平稳}。

\begin{definition}[严平稳]
随机过程 \(X(t)\) 称为严平稳，若对任意 \(n \ge 1\)，任意时刻 \(t_1, t_2, \dots, t_n \in T\) 以及任意时间平移 \(\tau\) 使得 \(t_i + \tau \in T\)，随机向量 \((X_{t_1}, \dots, X_{t_n})\) 与 \((X_{t_1+\tau}, \dots, X_{t_n+\tau})\) 具有相同的联合分布。
\end{definition}

严平稳的条件较强，在实际应用中往往难以验证。因此引入了宽平稳的概念，它仅要求一阶和二阶矩在时间平移下不变。

\begin{definition}[宽平稳]
随机过程 \(X(t)\) 称为宽平稳或二阶平稳，若满足：
\begin{enumerate}
    \item 均值函数为常数：\(m_{X}(t) = m_X\)，与 \(t\) 无关；
    \item 协方差函数仅依赖于时间差：\(C_{XX}(t, s) = \gamma(t-s)\)。
\end{enumerate}
\end{definition}

由于协方差函数和自相关函数只和时间差有关，所以我们也可以定义\(R_X(\tau)=R_{XX}(t,t+\tau)\)，它是良好定义的。

平稳性是时间序列分析、谱分析以及信号处理的核心假设，加上遍历性使得我们可以利用单条样本路径推断过程的统计特性，是我们做预测的基础。

\subsubsection{独立增量性}

独立增量性描述的是过程在不重合的时间区间上的变化量相互独立。这一性质使得过程具有无后效的累积特性，是构造许多基本随机过程的基础。

\begin{definition}[独立增量]
随机过程 \(X(t)\) 称为具有独立增量，若对任意 \(n \ge 2\) 和任意 \(0 \le t_1 < t_2 < \dots < t_n\)，增量
\[
X_{t_2} - X_{t_1},\; X_{t_3} - X_{t_2},\; \ldots,\; X_{t_n} - X_{t_{n-1}}
\]
相互独立。
\end{definition}

进一步地，若增量 \(X_{t+s} - X_s\) 的分布仅依赖于时间差 \(t\) 而与起点 \(s\) 无关，则称过程具有时齐性或者具有平稳独立增量。这类过程被称为Lévy过程，此时它是马尔可夫的，其代表性例子包括：

\begin{itemize}
    \item \textbf{泊松过程}：计数过程，增量服从泊松分布，用于描述稀有事件的发生；
    \item \textbf{维纳过程}：连续路径，增量服从正态分布，是金融数学与物理中的基本模型；
    \item \textbf{复合泊松过程}：用于建模带有随机跳幅的累积过程。
\end{itemize}

\subsubsection{高斯性}

高斯性是随机过程中另一类重要的性质。一个随机过程称为高斯过程，若其任意有限维分布均为多元正态分布。高斯过程在理论研究和实际应用中均占有特殊地位，这主要得益于中心极限定理。

\begin{definition}[高斯过程]
随机过程 \(X(t)\) 称为高斯过程（或称正态过程），若对任意 \(n \ge 1\) 和任意时刻 \(t_1, t_2, \dots, t_n \in T\)，随机向量 \((X_{t_1}, X_{t_2}, \dots, X_{t_n})\) 服从多元正态分布。
\end{definition}

高斯过程具有以下核心特征：

\begin{enumerate}
    \item \textbf{由均值函数与协方差函数唯一确定}：设均值函数 \(m_{X}(t)\)，协方差函数 \(C_{XX}(s, t) \)，则高斯过程的任意有限维分布完全由 \(m_{X}(t)\) 和 \(C_{XX}(s, t)\) 决定。这一性质使得高斯过程的统计推断变得相对简单——只需估计一阶和二阶矩即可刻画整个概率结构。

    \item \textbf{宽平稳与严平稳等价}：对于高斯过程，宽平稳（均值常数、协方差仅依赖于时间差）蕴涵严平稳。这是因为正态分布完全由其均值和协方差矩阵决定，若一、二阶矩具有平移不变性，则有限维分布亦具有平移不变性。

    \item \textbf{独立性等价于不相关性}：对于高斯过程，随机变量 \(X_s\) 与 \(X_t\) 独立当且仅当它们不相关（即协方差为零）。
\end{enumerate}